%% Список сокращений и условных обозначений по ГОСТ Р 7.0.11-2011, п. 5.4.
%% Перечень расположен столбцом: слева — сокращение, справа — расшифровка.
\chapter*{Список сокращений и условных обозначений}
\addcontentsline{toc}{chapter}{Список сокращений и условных обозначений}

\begin{longtable}{@{}p{0.20\textwidth}p{0.74\textwidth}@{}}
LLM     & большая языковая модель (Large Language Model) \\
MLLM    & мультимодальная большая языковая модель (Multimodal Large Language Model) \\
MoE     & смесь экспертов (Mixture of Experts) \\
RAG     & генерация с извлечением информации (Retrieval-Augmented Generation) \\
SFT     & дообучение с учителем (Supervised Fine-Tuning) \\
RLHF    & обучение с подкреплением на основе обратной связи от человека (Reinforcement Learning from Human Feedback) \\
RLAIF   & обучение с подкреплением на основе обратной связи от ИИ (Reinforcement Learning from AI Feedback) \\
DPO     & оптимизация прямых предпочтений (Direct Preference Optimization) \\
PPO     & проксимальная оптимизация политики (Proximal Policy Optimization) \\
GQA     & сгруппированное запросное внимание (Grouped-Query Attention) \\
MLA     & многоголовое латентное внимание (Multi-Head Latent Attention) \\
RoPE    & вращательное позиционное кодирование (Rotary Position Embedding) \\
ViT     & визуальный трансформер (Vision Transformer) \\
BM25    & функция ранжирования Best Matching 25 \\
RRF     & слияние по обратным рангам (Reciprocal Rank Fusion) \\
ANN     & приближённый поиск ближайших соседей (Approximate Nearest Neighbor) \\
InfoNCE & функция потерь Information Noise-Contrastive Estimation \\
BCE     & бинарная кросс-энтропия (Binary Cross-Entropy) \\
EMA     & экспоненциальное скользящее среднее (Exponential Moving Average) \\
OCR     & оптическое распознавание символов (Optical Character Recognition) \\
nDCG    & нормированный дисконтированный накопленный выигрыш (normalized Discounted Cumulative Gain) \\
\end{longtable}
